% !PRINTED 460
\weberpage{460}

\weq{3}{\Theta = \frac{\psi(\xi)}{F'(\xi)},}

and $\Theta$ is therefore contained in $\Omega(\xi)$. Since conversely every
quantity in $\Omega(\xi)$ is at the same time contained in $\Omega(\alpha, \beta, \gamma \dots)$, the
two fields are identical, i.\,e. we have

\[
\Omega(\xi) = \Omega(\alpha, \beta, \gamma \dots).
\]

But with this our assertion is proved.

\begin{center}
§. 144.

\gesperrt{Primitive and imprimitive fields.}
\end{center}

By the theorem last proved, we now restrict ourselves to
the consideration of algebraic fields $\Omega(\alpha)$, where $\alpha$
is the root of an equation in $\Omega$. Let this equation,
freed from multiple factors, be denoted by

\weq{1}{F(x) = 0}

and be of degree $m$. Its roots are as follows.

\weq{2}{\alpha, \alpha_1, \alpha_2 \dots \alpha_{m-1}.}

It may be that $F(x)$ is reducible in $\Omega$; then there will be one and
only one \gesperrt{irreducible factor} $f(x)$ of $F(x)$, so
that $\alpha$ is a root of the following equation.

\weq{3}{f(x) = 0}

Let the degree of $f(x)$ be $n$. The roots of (3), which are all
contained among the roots (2), are as follows.

\weq{4}{\alpha, \alpha_1, \alpha_2 \dots \alpha_{n-1}.}

The degree of the field $\Omega(\alpha)$ is then, by the stipulation
made in §.\,142, likewise $n$.

Each of the quantities (4) gives rise to an algebraic field,
and thus arise the fields

\weq{5}{\Omega(\alpha), \Omega(\alpha_1) \dots \Omega(\alpha_{n-1}),}

which we call the \gesperrt{fields conjugate to} $\Omega(\alpha)$. It may
be that these fields are all or in part identical, but they
may also all be different from one another, of which more
later.

By §.\,142 we obtain every quantity $\Theta$ in $\Omega(\alpha)$ if, in an
integral rational function $\chi(t)$ in $\Omega$, of degree at most
$(n-1)$, we put the number $\alpha$ for the variable $t$.
